\subsection{带标架的半单李代数的自同构}

回顾（第七章 §3 第 1 节），\(\operatorname{Aut}(\mathfrak{g})\) 表示 \(\mathfrak{g}\) 的自同构群。若 \(\mathfrak{h}\) 是 \(\mathfrak{g}\) 的一个嘉当子代数，我们用 \(\operatorname{Aut}(\mathfrak{g},\mathfrak{h})\) 表示 \(\mathfrak{g}\) 的使 \(\mathfrak{h}\) 稳定的自同构组成的群。假设 \(\mathfrak{h}\) 是分裂的，并设 R 是 \((\mathfrak{g},\mathfrak{h})\) 的根系。若 \(s\in\operatorname{Aut}(\mathfrak{g},\mathfrak{h})\) ，则 \(s|_{\mathfrak{h}}\) 的反步映射是 A(R) （R 的自同构群）的一个元素，我们在本节中将它记作 \(\varepsilon(s)\) 。于是

\[
\varepsilon : \operatorname{Aut} (\mathfrak {g}, \mathfrak {h}) \to \mathrm{A} (\mathrm{R})
\]

是一个群同态。

对任何根系 R 及 R 的任何基 B ，我们用 \(\operatorname{Aut}(R,B)\) 表示使 B 稳定的 R 的自同构组成的群。回顾（第六章 §1 第 5 节 命题 16 与 §4 第 2 节 命题 1 的推论），A(R) 是 \(\operatorname{Aut}(R,B)\) 与 W(R) 的半直积，且 A(R)/W(R) 典则同构于 R 的邓金图的自同构群。

命题 1. 设 \((\mathfrak{g},\mathfrak{h},\mathrm{B},(X_{\alpha})_{\alpha \in \mathrm{B}})\) 是一个带标架的半单李代数，R 是 \((\mathfrak{g},\mathfrak{h})\) 的根系。设 G 是这样的 \(s\in \operatorname {Aut}(\mathfrak{g},\mathfrak{h})\) 组成的集合：它们使 B 稳定，且使 \(s(X_{\alpha}) = X_{\varepsilon (s)\alpha}\) 对一切 \(\alpha \in \mathrm{B}\) 成立（换言之，即 \((\mathfrak{g},\mathfrak{h},\mathrm{B},(X_{\alpha})_{\alpha \in \mathrm{B}})\) 的自同构组成的集合）。则 \(\varepsilon\) 在 \(\mathrm{G}\) 上的限制是从 \(\mathrm{G}\) 到 \(\operatorname {Aut}(\mathrm{R},\mathrm{B})\) 的同构。

若 \(s \in G\) ，显然 \(\varepsilon(s) \in \operatorname{Aut}(\mathrm{R}, \mathrm{B})\) 。另一方面，映射

\[
\varepsilon|_{\mathrm{G}}: \mathrm{G} \rightarrow \operatorname{Aut} (\mathrm{R}, \mathrm{B})
\]

由 §4 第 4 节 定理 2 是双射的。

